\documentclass[12pt]{article}

% ==================== PAQUETES ====================
\usepackage[utf8]{inputenc}
\usepackage[T1]{fontenc}
\usepackage[spanish]{babel}
\usepackage{amsmath}
\usepackage{amsfonts}
\usepackage{amssymb}
\usepackage{geometry}
\usepackage{xcolor}
\usepackage{enumitem}

\geometry{margin=2.5cm}
\setlength{\parindent}{0pt}
\setlength{\parskip}{0.6em}

% ==================== METADATOS ====================
\title{Manual introductorio de IA y Visión en la Agroindustria}
\author{John Jairo Leal Gómez}
\date{Semanas 05 y 06}

\begin{document}

\maketitle
\tableofcontents
\newpage

% ==================================================
% SEMANA 05
% ==================================================
\part{Semana 05: Clasificación y lógica industrial}

\section{Incertidumbre y selección: entropía e información mutua}

\subsection{Explicación conceptual}
En problemas de clasificación, una de las preguntas más importantes es: ``\textit{Qué tan difícil es decidir a qué clase pertenece un objeto?}''.
La \textbf{entropía} mide precisamente ese nivel de incertidumbre.
Si todos los elementos pertenecen a una sola clase, la incertidumbre es baja.
Si las clases están mezcladas de manera equilibrada, la incertidumbre es alta.

En agroindustria, esto aparece cuando se desea clasificar frutas, hojas o semillas en categorías como ``sano'' y ``defectuoso''.
Si antes de observar una característica no sabemos bien a qué clase pertenece un elemento, tenemos mucha incertidumbre.
Si una característica, por ejemplo una mancha o un olor, reduce esa duda, entonces esa característica aporta información útil.

La \textbf{información mutua} mide cuánto ayuda una variable a reducir la incertidumbre sobre otra.
En términos prácticos, nos dice qué tan valiosa es una característica para clasificar correctamente.

\subsection{Explicación matemática}
Sea \(Y\) una variable aleatoria que representa la clase de un producto.
La entropía de \(Y\) se define como:
\[
H(Y) = - \sum_{y \in \mathcal{Y}} p(y)\log_2 p(y)
\]

donde \(p(y)\) es la probabilidad de la clase \(y\).

Si usamos una característica \(X\), la incertidumbre restante sobre \(Y\) después de observar \(X\) se mide con la entropía condicional:
\[
H(Y \mid X)
\]

La información mutua entre \(X\) y \(Y\) se define como:
\[
I(X;Y) = H(Y) - H(Y \mid X)
\]

Interpretación:
\begin{itemize}
    \item Si \(I(X;Y) = 0\), la característica no ayuda a clasificar.
    \item Si \(I(X;Y)\) es grande, la característica reduce bastante la incertidumbre.
    \item Si \(I(X;Y) = H(Y)\), entonces la característica permite identificar perfectamente la clase.
\end{itemize}

\subsection{Ejemplo}
Supongamos un lote de \(10\) mangos:
\begin{itemize}
    \item \(6\) están sanos.
    \item \(4\) están defectuosos.
\end{itemize}

La entropía inicial de la clase es:
\[
H(Y) = -\left(\frac{6}{10}\log_2\frac{6}{10} + \frac{4}{10}\log_2\frac{4}{10}\right)
\]

Calculando:
\[
H(Y) \approx 0.971
\]

Ahora usamos la característica ``mancha oscura'':
\begin{itemize}
    \item \(5\) mangos tienen mancha; de ellos, \(4\) son defectuosos y \(1\) es sano.
    \item \(5\) mangos no tienen mancha; de ellos, los \(5\) son sanos.
\end{itemize}

La entropía del grupo con mancha es:
\[
H(Y \mid X=\text{mancha}) =
-\left(\frac{4}{5}\log_2\frac{4}{5} + \frac{1}{5}\log_2\frac{1}{5}\right)
\approx 0.722
\]

La entropía del grupo sin mancha es:
\[
H(Y \mid X=\text{sin mancha}) =
-\left(\frac{5}{5}\log_2\frac{5}{5}\right) = 0
\]

Entonces, la entropía condicional total es:
\[
H(Y \mid X) = \frac{5}{10}(0.722) + \frac{5}{10}(0) = 0.361
\]

Finalmente, la información mutua es:
\[
I(X;Y) = 0.971 - 0.361 = 0.610
\]

Esto significa que la característica ``mancha oscura'' reduce bastante la incertidumbre y es útil para clasificar los mangos.

\subsection{Ejercicios}
\begin{enumerate}[label=\arabic*.]
    \item En un lote de \(10\) frutas, \(8\) son sanas y \(2\) están dañadas. Calcule la entropía de la clase.
    \item Una característica divide un lote en dos grupos: en el primero hay \(4\) sanos y \(1\) dañado; en el segundo hay \(4\) sanos y \(1\) dañado. \`La información mutua será alta o baja? Justifique.
    \item Explique con sus palabras por qué una característica que separa perfectamente las clases tiene mayor utilidad.
    \item Si \(H(Y)=1\) y después de observar \(X\) se obtiene \(H(Y \mid X)=0.2\), calcule \(I(X;Y)\).
\end{enumerate}

% --------------------------------------------------

\section{Árboles de decisión y métricas de calidad}

\subsection{Explicación conceptual}
Un árbol de decisión es un modelo que clasifica mediante una secuencia de preguntas.
Cada nodo del árbol formula una condición, por ejemplo: ``\textit{¿Tiene mancha?}'', ``\textit{¿Huele fermentado?}'' o ``\textit{¿La textura es blanda?}''.
Dependiendo de la respuesta, el elemento sigue una rama hasta llegar a una decisión final.

La idea central es elegir, en cada paso, la pregunta que mejor separe las clases.
Por eso los árboles de decisión están muy relacionados con la entropía y la información mutua.

Sin embargo, construir un árbol no es suficiente.
Después de entrenarlo, debemos evaluar qué tan bueno es.
En contextos industriales, no basta con medir solo el porcentaje total de aciertos.
A veces es más importante detectar todos los casos defectuosos que maximizar el número global de aciertos.

\subsection{Explicación matemática}
La medida más usada para elegir una división es la \textbf{ganancia de información}:
\[
IG(S,A) = H(S) - \sum_{i=1}^{k} w_i H(S_i)
\]

donde:
\begin{itemize}
    \item \(S\) es el conjunto original.
    \item \(A\) es el atributo usado para dividir.
    \item \(S_i\) son los subconjuntos generados.
    \item \(w_i = \frac{|S_i|}{|S|}\) es el peso de cada subconjunto.
\end{itemize}

Entre mayor sea \(IG(S,A)\), mejor es la pregunta para dividir el conjunto.

Para evaluar el clasificador, se usan las siguientes cantidades:
\begin{itemize}
    \item \(TP\): verdaderos positivos.
    \item \(FP\): falsos positivos.
    \item \(FN\): falsos negativos.
    \item \(TN\): verdaderos negativos.
\end{itemize}

A partir de ellas se definen las métricas:
\[
\text{Precision} = \frac{TP}{TP+FP}
\]

\[
\text{Recall} = \frac{TP}{TP+FN}
\]

\[
F1 = 2 \cdot \frac{\text{Precision} \cdot \text{Recall}}{\text{Precision} + \text{Recall}}
\]

\[
\text{Accuracy} = \frac{TP+TN}{TP+TN+FP+FN}
\]

Interpretación:
\begin{itemize}
    \item \textbf{Precision}: de todo lo que el modelo marcó como positivo, cuánto era realmente positivo.
    \item \textbf{Recall}: de todos los positivos reales, cuántos detectó el modelo.
    \item \textbf{F1}: equilibrio entre precision y recall.
    \item \textbf{Accuracy}: proporción total de aciertos.
\end{itemize}

\subsection{Ejemplo}
Supongamos que un sistema evalúa \(100\) frutas:
\begin{itemize}
    \item \(20\) estaban realmente defectuosas.
    \item El modelo detectó bien \(18\) defectuosas.
    \item Marcó por error \(6\) frutas sanas como defectuosas.
    \item Dejó pasar \(2\) defectuosas como si fueran sanas.
\end{itemize}

Entonces:
\[
TP = 18,\quad FP = 6,\quad FN = 2,\quad TN = 74
\]

La \textit{precision} es:
\[
\text{Precision} = \frac{18}{18+6} = \frac{18}{24} = 0.75
\]

El \textit{recall} es:
\[
\text{Recall} = \frac{18}{18+2} = \frac{18}{20} = 0.90
\]

El puntaje \(F1\) es:
\[
F1 = 2\cdot\frac{0.75\cdot 0.90}{0.75+0.90}
\approx 0.818
\]

La exactitud global es:
\[
\text{Accuracy} = \frac{18+74}{100} = 0.92
\]

Aunque la exactitud es alta, el dato más importante en este caso puede ser el \textit{recall}, porque dejar pasar fruta defectuosa puede afectar todo el lote.

\subsection{Ejercicios}
\begin{enumerate}[label=\arabic*.]
    \item Un modelo tiene \(TP=30\), \(FP=10\), \(FN=5\) y \(TN=55\). Calcule \textit{precision}, \textit{recall}, \(F1\) y \textit{accuracy}.
    \item Explique por qué, en control de calidad, un falso negativo puede ser más costoso que un falso positivo.
    \item Si un árbol puede dividir un conjunto usando dos atributos, uno con \(IG=0.12\) y otro con \(IG=0.35\), \`cuál debería elegir primero y por qué?
    \item Un modelo tiene \textit{recall} del \(100\%\) pero \textit{precision} del \(10\%\). Describa qué ocurriría en la operación real de una planta agroindustrial.
\end{enumerate}

% ==================================================
% SEMANA 06
% ==================================================
\part{Semana 06: CNN y visión artificial}

\section{Tensores y convoluciones}

\subsection{Explicación conceptual}
Para un sistema de inteligencia artificial, una imagen no es una ``foto'' como la ve una persona, sino un conjunto ordenado de números.
Ese conjunto se representa como un \textbf{tensor}.

Si la imagen está en escala de grises, puede representarse como una matriz de tamaño \(alto \times ancho\).
Si está en color, normalmente se representa como un tensor de tamaño \(alto \times ancho \times canales\), donde los canales suelen ser rojo, verde y azul.

Las redes convolucionales usan \textbf{filtros} o \textbf{kernels}, que son matrices pequeñas que recorren la imagen.
Cada filtro busca un patrón específico, como bordes, manchas, texturas o cambios bruscos de color.
En agricultura, esto es útil para detectar síntomas de plagas, enfermedades o defectos superficiales.

\subsection{Explicación matemática}
Sea una imagen \(X\) y un filtro \(K\).
La operación de convolución discreta en \(2\) dimensiones puede escribirse como:
\[
Y(i,j) = \sum_{m}\sum_{n} X(i+m,j+n)\,K(m,n)
\]

donde:
\begin{itemize}
    \item \(X\) es la imagen de entrada.
    \item \(K\) es el filtro.
    \item \(Y\) es la salida resultante.
\end{itemize}

En muchas implementaciones de redes convolucionales, esta operación corresponde en realidad a una correlación discreta, pero en cursos introductorios suele llamarse igualmente ``convolución''.

Si una imagen en escala de grises tiene tamaño \(H \times W\), entonces contiene:
\[
H \cdot W
\]
números.

Si la imagen es RGB, contiene:
\[
H \cdot W \cdot 3
\]
valores.

Sin relleno y con paso \(1\), una imagen de tamaño \(H \times W\) filtrada con un kernel de tamaño \(h \times w\) produce una salida de tamaño:
\[
(H-h+1)\times(W-w+1)
\]

\subsection{Ejemplo}
Considere la siguiente imagen en escala de grises:
\[
X=
\begin{bmatrix}
1 & 2 & 0 \\
0 & 1 & 3 \\
2 & 1 & 0
\end{bmatrix}
\]

y el siguiente filtro:
\[
K=
\begin{bmatrix}
1 & 0 \\
0 & -1
\end{bmatrix}
\]

Aplicando el filtro con paso \(1\) y sin relleno, la salida tendrá tamaño:
\[
(3-2+1)\times(3-2+1)=2\times 2
\]

Ahora calculamos cada posición.

Primera posición:
\[
Y(1,1)=
1\cdot 1 + 2\cdot 0 + 0\cdot 0 + 1\cdot(-1)=0
\]

Segunda posición:
\[
Y(1,2)=
2\cdot 1 + 0\cdot 0 + 1\cdot 0 + 3\cdot(-1)=-1
\]

Tercera posición:
\[
Y(2,1)=
0\cdot 1 + 1\cdot 0 + 2\cdot 0 + 1\cdot(-1)=-1
\]

Cuarta posición:
\[
Y(2,2)=
1\cdot 1 + 3\cdot 0 + 1\cdot 0 + 0\cdot(-1)=1
\]

Por tanto, la salida es:
\[
Y=
\begin{bmatrix}
0 & -1 \\
-1 & 1
\end{bmatrix}
\]

\subsection{Ejercicios}
\begin{enumerate}[label=\arabic*.]
    \item Si una imagen en escala de grises tiene tamaño \(10\times 10\), \`cuántos valores contiene el tensor?
    \item Si una imagen RGB tiene tamaño \(10\times 10\), \`cuántos valores contiene el tensor?
    \item Si una imagen tiene tamaño \(28\times 28\) y se aplica un filtro de \(3\times 3\) con paso \(1\) y sin relleno, \`cuál será el tamaño de la salida?
    \item Aplique el filtro
    \[
    \begin{bmatrix}
    1 & 0 \\
    0 & -1
    \end{bmatrix}
    \]
    sobre la matriz
    \[
    \begin{bmatrix}
    2 & 1 & 0 \\
    1 & 2 & 1 \\
    0 & 1 & 2
    \end{bmatrix}
    \]
    y calcule la salida completa.
\end{enumerate}

% --------------------------------------------------

\section{CNN en agroindustria: detección de plagas y enfermedades}

\subsection{Explicación conceptual}
Una red neuronal convolucional, o CNN, procesa imágenes en varias etapas.
Primero identifica patrones simples, como bordes o manchas.
Luego combina esos patrones para reconocer estructuras más complejas, como zonas enfermas, perforaciones, decoloraciones o texturas anormales.

En una aplicación agroindustrial, una CNN puede analizar una imagen de una hoja y decidir si está sana o si presenta síntomas de una enfermedad.
Esto permite apoyar tareas de monitoreo, selección y diagnóstico temprano.

Una arquitectura básica suele incluir:
\begin{itemize}
    \item \textbf{Capas convolucionales}, que extraen características.
    \item \textbf{Funciones de activación}, que introducen no linealidad.
    \item \textbf{Capas de \textit{pooling}}, que reducen tamaño y conservan rasgos importantes.
    \item \textbf{Capas finales de clasificación}, que emiten la decisión.
\end{itemize}

\subsection{Explicación matemática}
Después de una convolución, suele aplicarse una función de activación como ReLU:
\[
\text{ReLU}(z)=\max(0,z)
\]

Luego puede aplicarse una operación de \textit{max-pooling} sobre una región \(R\):
\[
P = \max_{(i,j)\in R} X(i,j)
\]

Si el problema es binario, por ejemplo ``hoja sana'' frente a ``roya'', la salida final puede modelarse con una función sigmoide:
\[
\hat{y} = \sigma(z)=\frac{1}{1+e^{-z}}
\]

donde \(\hat{y}\) es la probabilidad estimada de que la imagen pertenezca a la clase positiva.

La regla de decisión más simple es:
\begin{itemize}
    \item Si \(\hat{y}\geq 0.5\), clasificar como ``enferma''.
    \item Si \(\hat{y}<0.5\), clasificar como ``sana''.
\end{itemize}

\subsection{Ejemplo}
Supongamos que una CNN analiza una imagen de una hoja de café y en la última capa obtiene:
\[
z=2.2
\]

Entonces, la probabilidad estimada de ``roya'' es:
\[
\hat{y}=\frac{1}{1+e^{-2.2}} \approx 0.900
\]

Como \(\hat{y}\geq 0.5\), el sistema clasifica la hoja como:
\[
\text{``Roya''}
\]

Interpretación práctica:
\begin{itemize}
    \item La red encontró patrones compatibles con la enfermedad.
    \item El valor \(0.900\) indica alta confianza relativa en la clase positiva.
    \item En campo o planta, esta salida puede usarse para activar una alerta o enviar la muestra a revisión.
\end{itemize}

\subsection{Ejercicios}
\begin{enumerate}[label=\arabic*.]
    \item Si una red produce \(\hat{y}=0.23\), \`cómo se clasificaría la hoja usando un umbral de \(0.5\)?
    \item Calcule \(\hat{y}\) cuando \(z=0\). Interprete el resultado.
    \item Explique por qué una capa de \textit{pooling} puede ayudar a reducir el costo computacional.
    \item En un sistema de detección de plagas, \`por qué puede ser riesgoso tener un \textit{recall} muy bajo?
\end{enumerate}

% --------------------------------------------------

\section{Ejercicios integradores}

\begin{enumerate}[label=\arabic*.]
    \item Un modelo de detección de plagas tiene \textit{recall} del \(100\%\) y \textit{precision} del \(10\%\). Explique qué le ocurre al agricultor en la práctica.
    \item Una imagen de tamaño \(10\times 10\) en escala de grises entra a una CNN. \`Cuántos valores tiene el tensor de entrada?
    \item Una base de datos tiene dos clases equilibradas: ``sano'' y ``enfermo''. \`La entropía inicial será alta o baja? Justifique.
    \item En una planta de selección, compare dos modelos: uno con \textit{accuracy} del \(95\%\) y \textit{recall} del \(60\%\), y otro con \textit{accuracy} del \(90\%\) y \textit{recall} del \(95\%\). \`Cuál elegiría para evitar que salgan productos defectuosos y por qué?
\end{enumerate}

\end{document}
